\printedpage{54}
% head: A. Grothendieck : Dérivateurs, Ch. VI, Édition des Universités de Montpellier II et Paris VII

\begin{flushright}
16.11.90
\end{flushright}

I am going to try all the same to do some ``roughing-out'' work, for the diagram of the
15 categories looks a bit bushy, % [sic] peut -> peu, adverb "un peu" intended
after all! I give myself
\[
(\cM,W) \qquad W\subseteq \Fl\,\cM,
\]
an $S\in\Ob\,\cM$, whence $\cM/S$, with $W_S\subseteq\Fl(\cM/S)$ induced by $W$. We are going to put
\[
\Ho^{\mathrm{lc}}(S)=(\cM/S)W_S^{-1},
\]
this will play the role of relatively constant homotopy types. I suppose on the other hand
that we have singled out a class $F\subseteq\Fl(\cM)$ of arrows, playing the role of ``fibrations'' (these
will be the Hot-fibrant morphisms in Cat, giving rise to the subcategory $\Fibhot S$
of $\mathrm{Cat}/S$). I foresee that we will have to
\msp{page 86}suppose, besides $\cM$ being stable under finite $\varprojlim$, that $F$ is stable under arbitrary base change,
and stable under composition, that it contains the isomorphisms. And above all that we
have a factorization theorem
\[
\begin{tikzcd}
X \arrow[rr, "f"] \arrow[dr, "i_X"'] & & S \\
& \bar{X} \arrow[ur, "\bar{f}"'] &
\end{tikzcd}
\qquad\qquad\qquad
i_X\in W,\ \bar{f}\in F
\]
for % [sic] pout -> pour, misprint
every arrow $f$ in $\cM$. With \textsc{Quillen} one even requires that $i_X$ have the right lifting
property with respect to $F$, but that should not be essential, since I did not
have to use this fact in Cat. When I work in the category $\cM/S$, where $f$ is
understood for an object $X$ of the said category, I will write (by abuse of language) $C(X)$ for $\bar{X}$,
although it is not said that we have a canonical choice of $C(X)$, still less that $C(X)$
depends functorially on $X$.

Let us denote by
\[
\Fib S\subseteq\cM/S \tag*{$(*)$}
\]
the subcategory of $\cM/S$ formed by the $X$ whose structural morphism $f_X:X\longrightarrow S$ is
in $F$. Let us denote (by abuse of notation) by $W_S$ the localizer induced on $\Fib S$ by $W_S$,
then the inclusion $(*)$ defines
\[
(\Fib S)W_S^{-1}\xrightarrow{\ \delta\ }(\cM/S)W_S^{-1}\overset{\text{def}}{=}\Ho^{\mathrm{lc}}_W(S).
\]

\msp{page 87}A first question is to see that this is an equivalence. We see that it is essentially
surjective thanks to the existence of the $C(X)$, $i_X$. That is not much! We are going to try to define
an arrow in the opposite direction
\[
\bar{C}:(\cM/S)W_S^{-1}\longrightarrow(\Fib S)W_S^{-1},
\]
% foot: Version 9 novembre 2001
\printedpage{55}
% head: A. Grothendieck : Dérivateurs, Ch. VI, Édition des Universités de Montpellier II et Paris VII
i.e. a \textit{functor}
\[
\gamma:\cM/S\longrightarrow(\Fib S)W_S^{-1}
\]
which transforms $W_S\subseteq\Fl(\cM/S)$ into isomorphisms. We will choose, for $X$ in $\cM/S$, a
$C(X)$ in $\Fib S$ and an
\[
i_X:X\longrightarrow C(X)\ , \qquad i_X\in W_S,
\]
and we will put
\[
\gamma(X)=C(X).
\]
The difficulty is to show that this does not depend on the choices made, up to canonical
isomorphism, and that it is functorial in $X$ - and the second implies the first.

To see that it works, clearly we have to suppose something more. I see two
hypotheses which, either of them, could make it work: one is that we already give ourselves
$C$ as a functor and $i_X$ functorial in $X$ (this was the case in Cat just now).
The other is that
\msp{page 88}the $i_X$ have, with respect to $F$, the lifting property in the Quillen style. % [sic] à là -> à la, accent on article "la"

I begin with the hypothesis that $C$ is already a functor. The diagram
\[
\begin{tikzcd}[row sep=large]
X \arrow[r, "i_X"] \arrow[d, "u"'] & CX \arrow[d, "C(u)"] \\
X' \arrow[r, "i_{X'}"] & CX',
\end{tikzcd}
\]
where $i_X,\ i_{X'}\in W$, shows that $C$ transforms $W_S$ into $W_S$, hence into isomorphisms of
$\Ho^{\mathrm{lc}}_W(S)$. So we have got the sought-for functor $\gamma$. In fact we have functors
\[
\begin{tikzcd}
\Fib S \arrow[r, hook, "\overset{\text{incl.}}{D}"] & \cM/S \arrow[l, bend left=50, "C"]
\end{tikzcd}
\]
compatible with the localizers $W_S$, inducing $\delta$ and $\gamma$, and we have
\[
\mathrm{id}_{\cM/S}\longrightarrow DC, \qquad \text{in $W_S$ argument by argument,}
\]
for we have
\[
X\xrightarrow{\ i_X\ }DC(X)=C(X)
\]
with $i_X\in W_S$. Likewise, we have
\[
\mathrm{id}_{\Fib S}\longrightarrow CD \qquad \text{in $W_S$ argument by argument,}
\]
% foot: Version 9 novembre 2001
\printedpage{56}
% head: A. Grothendieck : Dérivateurs, Ch. VI, Édition des Universités de Montpellier II et Paris VII
with the same
\[
X\xrightarrow{\ i_X\ }CX \qquad \text{(this time $X$ in $\Fib S$).}
\]

\msp{page 89}So we have without tears the two functors $\delta$, $\gamma$ as quasi-inverses of each other. And the
lemma on page 11 even shows that for \textit{every} full subcategory $\cF$ of $\cM/S$, such that
\[
\Fib S\subseteq\cF\subseteq\cM/S,
\]
the inclusions induce equivalences on the $W_S$-localized categories
\[
(\Fib S)W_S^{-1}\xrightarrow{\ \sim\ }\cF W_S^{-1}\xrightarrow{\ \sim\ }\underbrace{(\cM/S)W_S^{-1}}_{\overset{\text{def}}{=}\Ho^{\mathrm{lc}}_W(S)}.
\]
I am now going to treat the case of Quillen type, to convince myself that the existence of the
functor $C$, and the particular choice one can make of it, is irrelevant. (This last point
is moreover obvious a posteriori, when we know that $\delta$ is an equivalence, for the quasi-inverse
$\gamma$ will then be unique up to unique isomorphism, and the choice of a functor $C$, and
of a functorial homomorphism $i:\mathrm{id}\longrightarrow C$, is only one way among others of making
the quasi-inverse explicit.)

\msp{page 90}We therefore place ourselves under the Quillenian hypothesis on the $i_X$. We choose $C(X)$, $i_X$ for
each $X$ in $\cM/S$, and we put as before
\[
\gamma(X)=C(X) \qquad \text{as an object of $(\cM/S)W_S^{-1}$.}
\]
Let us show that this is functorial in $X$. Let
\[
u:X\longrightarrow X'
\]
be in $\cM/S$, let us consider the harmless (commutative) diagram
\[
\begin{tikzcd}[row sep=large]
X \arrow[rr, "u"] \arrow[d, "i_X"'] \arrow[drr, loosely dotted, "v=i_{X'}u"] & & X' \arrow[d, "i_{X'}"] \\
C(X) \arrow[rr, dotted, "\bar{u}"] \arrow[dr] & & C(X') \arrow[dl, "\in F"] \\
& S, &
\end{tikzcd}
\]
let us complete it with the coarsely dotted arrow, $v=i_{X'}u$ (making the upper triangle
commutative). (\footnote[54]{Lemma: let
% footnote continues on printed p. 57
\[
\begin{tikzcd}[ampersand replacement=\&, row sep=large]
X \arrow[r, "u"] \arrow[d, "i"'] \& X' \arrow[d] \\
C \arrow[r, dotted] \& C'
\end{tikzcd}
\]
be in $\cM/S$ with $i\in\mathrm{Cof}_0$, $(C'\longrightarrow S)\in F=\mathrm{Fib}$, then the dotted arrow exists in $\cM/S$, so as
to make the square commutative. \texttt{[cf. Lemma p. 91]}}). The hypothesis on $i_X$ implies that the finely dotted arrow $\bar{u}$ exists,
so as to have two more commutative triangles. But this means exactly that there is
an $S$-\textit{morphism} $\bar{u}$, i.e. an arrow in $\cM/S$, making the square commutative
% foot: Version 9 novembre 2001
\printedpage{57}
% head: A. Grothendieck : Dérivateurs, Ch. VI, Édition des Universités de Montpellier II et Paris VII
\[
\begin{tikzcd}[row sep=large]
X \arrow[r, "u"] \arrow[d, "i_X"'] & X' \arrow[d, "i_{X'}"] \\
C(X) \arrow[r, "\bar{u}"] & C(X'),
\end{tikzcd}
\]
\msp{page 91}and we must take
\[
\gamma(u)=\ \text{image of $\bar{u}$ in $\Fl((\Fib S)W_S^{-1})$.}
\]
We have to prove that this arrow does not depend on the choice of $\bar{u}$. We must isolate a Lemma,
where we denote by $\mathrm{Cof}_0$ the set of arrows of $\cM$ which have % [sic] qu'ont -> qui ont, relative pronoun
the lifting property with
respect to the elements of $\mathrm{Fib}=F$. ($\mathrm{Cof}_0$ plays the role of trivial cofibrations.)

\textbf{Lemma.} \textit{Let there be a solid-line diagram in $\cM/S$}
\[
\begin{tikzcd}[row sep=large]
X \arrow[r, "u"] \arrow[d, "i"'] & X' \arrow[d, "i'"] \\
C \arrow[r, dotted, "v"] & C'
\end{tikzcd}
\]
\textit{with $i\in$ Cof$_0$, and $C'\in$ Ob$\,\Fib S$ i.e. $(C'\longrightarrow S)\in$ Fib. Then $v$ exists, so as to make
the square commutative. Moreover, if $i\in W$ and $C\in\Ob\,\Fib S$, the class of $v$ as an arrow of
$(\Fib S)W_S^{-1}$ is independent of the choice made.}

It remains to prove this independence,
\msp{page 92}which will follow at once from the corollary:

\textbf{Corollary.} \textit{Let there be a diagram in $\cM/S$}
\[
\begin{tikzcd}
X \arrow[r, "i"] & C \arrow[r, shift left, "v_1"] \arrow[r, shift right, "v_2"'] & C',
\end{tikzcd}
\]
\textit{with $C$ and $C'$ in $\Fib S$ (i.e. $C\longrightarrow S$ and $C'\longrightarrow S$ in Fib), and $i\in$ Cof$_0$ and $i\in W$. If
$v_1i=v_2i$, then $v_1$ and $v_2$ define the same arrow of $(\Fib S)W_S^{-1}$.}
% foot: Version 9 novembre 2001

\printedpage{58}
% head: A. Grothendieck : Dérivateurs, Ch. VI, Édition des Universités de Montpellier II et Paris VII
\textsc{Proof}. Let
\[
v=(v_1,v_2):C\longrightarrow C'\times C',
\]
\texttt{[rather} $C'\times_S C'$\texttt{]} which by hypothesis makes the diagram commutative
\[
\begin{tikzcd}[row sep=large]
X \arrow[r, "w"] \arrow[d, "i"'] & C' \arrow[d, "\mathrm{diag}_{C'}"] \\
C \arrow[r, "v"] & C'\times_S C',
\end{tikzcd}
\]
where $w=v_1i=v_2i$. Let us factor $\mathrm{diag}_{C'}$ as
\[
C'\xrightarrow{\ i'\ }\bar{C}'\xrightarrow{\ \delta'\ }C'\times_S C', \qquad \text{with $\delta'\in\mathrm{Fib}$,}
\]
and $i'\in W$. We find a commutative square
\[
\begin{tikzcd}[row sep=large]
X \arrow[r, "w'=i'w"] \arrow[d, "i"'] & \bar{C}' \arrow[d, "\delta'"] & \\
C \arrow[r, "v"] \arrow[ur, dotted, "\bar{v}"] & C'\times_S C' \arrow[r, shift left, "\mathrm{pr}_1"] \arrow[r, shift right, "\mathrm{pr}_2"'] & C'
\end{tikzcd}
\]
with $i\in\mathrm{Cof}_0$, $\delta'\in\mathrm{Fib}$, whence the existence of the dotted arrow, making the two triangles
commutative. We want to prove $\mathrm{pr}_1\circ v=\mathrm{pr}_2\circ v$ in $(\Fib S)W_S^{-1}$, as $v=\delta'\bar{v}$ it suffices
to prove that we have $\mathrm{pr}_1\delta'=\mathrm{pr}_2\delta'$ in $(\Fib S)W_S^{-1}$.

\msp{page 93}We only have to check first (to justify the argument) that this assertion makes sense,
i.e. that $\bar{C}'$ is itself also in $\Fib S$. But the hypotheses of stability of $F=\mathrm{Fib}$ under
base change and composition ensure that $1^o$) $C'\times_S C'$ is in $\Fib S$ (which is
stable under products), and consequently, $2^o$) that $\bar{C}'$ is also, since $\delta'\in\mathrm{Fib}$. As $C'$ is
also in $\Fib S$ and $C'\xrightarrow{\ i'\ }\bar{C}'$ is in $W_S$, to check $\mathrm{pr}_1\delta'=\mathrm{pr}_2\delta'$ in the localized category of
$\Fib S$, it suffices to see that $(\mathrm{pr}_1\delta')i'=(\mathrm{pr}_2\delta')i'$, now since $\delta'i'=\mathrm{diag}_{C'}$, we have the equality even
in $\Fib S$ (without localizing).

We have therefore defined $\gamma(X)$, $\gamma(u)$ for objects and arrows of $\cM/S$. Transitivity for $\gamma(u)$ is
clear, and also $\gamma(\mathrm{id})=\mathrm{id}$. So we have got the functor
\[
\gamma:\cM/S\longrightarrow(\Fib S)W_S^{-1}.
\]
Let now $\cF$ be a full subcategory of $\cM/S$, with
\[
\Fib S\subseteq\cF\subseteq\cM/S,
\]
let us consider $\gamma$ as a functor with values in $\cF W_S^{-1}$. (I forgot to check that $\gamma$ transforms
$W_S$ into $W_S$, but this is tautological.) Thus we have functors
% foot: Version 9 novembre 2001